% TODO: 論文題目等の情報を以下に記入

\newcommand{\jtitle}{論文タイトルは４６文字で２行に収まります１２３４５６７８９０１２３４５６７８９０１２３４５６}  % 卒業論文の題名（日）
\newcommand{\etitle}{Title in English Within Two Lines: Lorem Ipsum Dolor Sit Amet, Consectetur Adipiscing Elit, Sed Do}   % 論文題目（英）
%タイトル：APAスタイルとします．特に大文字化の指示がない限り、単語はすべて小文字で表記します。例えば、文頭の単語は大文字で表記しますが、文頭の人物名が小文字で始まる場合は除きます。
%参考：Capitalization https://apastyle.apa.org/style-grammar-guidelines/capitalization
%参考：Title Case Converter – A Smart Title Capitalization Tool https://titlecaseconverter.com/
\newcommand{\jauthor}{姓 名}      % 著者名（日）
\newcommand{\eauthor}{LASTNAME Firstname} % 著者名（英）
\newcommand{\jadvisor}{姓 名}   % 指導教員名（日）
\newcommand{\eadvisor}{LASTNAME Firstname}  % 指導教員名（英）
%教員名の表記は以下のサイトの表記に従うこと
%教員 – 公立はこだて未来大学 -Future University Hakodate- https://www.fun.ac.jp/faculty
\newcommand{\jdate}{20XX年1月XX日}  % 論文提出日   (日)
\newcommand{\edate}{January XXth, 20XX}  % 論文提出年月　(英)
\newcommand{\jkeywords}{キーワード１, キーワード2, キーワード3} % キーワード（日）
\newcommand{\ekeywords}{Keyword1, Keyword2, Keyword3}   % キーワード（英）
\newcommand{\eshorttitle}{Your Short English Title Here}    % 短縮英題題名（おおよそ8 words以内）
\newcommand{\jdepartment}{情報アーキテクチャ学科}    % 学科名（日）
%\newcommand{\jdepartment}{複雑系知能学科}    % 学科名（日）
\newcommand{\jcourse}{情報システムコース}    % コース名（日）
%\newcommand{\jcourse}{高度ICTコース}    % コース名（日）
%\newcommand{\jcourse}{情報デザインコース}    % コース名（日）
%\newcommand{\jcourse}{複雑系コース}    % コース名（日）
%\newcommand{\jcourse}{知能システムコース}    % コース名（日）
\newcommand{\studentID}{1099999}    % 学籍番号
\newcommand{\edepartment}{Department of Media Architecture}    % 学科名（英）
%\newcommand{\edepartment}{Department of Complex and Intelligent Systems}    % 学科名（英）
\newcommand{\ecourse}{Information Systems Course}    % コース名（英）
%\newcommand{\ecourse}{Advanced ICT Course}    % コース名（英）
%\newcommand{\ecourse}{Information Design Course}    % コース名（英）
%\newcommand{\ecourse}{Complex Systems Course}    % コース名（英）
%\newcommand{\ecourse}{Intelligent Systems Course}    % コース名（英）
